% (User input window)

\begin{remark}[编译体积与可见维度的分离：对称幂商核的多项式体积上界与最小 realization 的严格压缩]\label{rem:pom-compile-volume-vs-visible-dimension}
对 \(q\)-重碰撞，直接对底层核做 \(q\) 份同步直积会产生指数态空间 \(|Q|^q\)（底层状态数为 \(|Q|\)）。
在具有叶子置换对称的情形下（注 \ref{rem:pom-star-sympress}），对称群取商把指数态空间压缩为占据数空间 \(N_q\)，其基数满足多项式上界
\[
|N_q|\ \le\ \binom{q+|Q|-1}{|Q|-1}.
\]
相应地，对称幂商核 \(\widetilde A_q\) 给出一个可执行的“编译后端体积”量：\(\log|N_q|\) 记录的是组合态空间大小，而非谱语义的真实自由度。
\par
另一方面，序列侧的 Hankel 最小化（注 \ref{rem:pom-minh-functor}；注 \ref{rem:pom-sympower-then-minh}）把同一碰撞矩序列 \(S_q(m)\) 压缩到其最小整系数递推阶数 \(d_q\)，等价于最小 realization 维数（命题 \ref{prop:pom-moment-minrecurrence-hankel}）。
在已核验区间内，\(d_q\) 往往显著小于上述编译体积（例如 \(d_4=5\) 见表 \ref{tab:fold_collision_moment_spectrum_k2_8}；而更高阶出现共振缺口 \(\Delta(q)=(2\lfloor q/2\rfloor+1)-d_q>0\)，见注 \ref{rem:pom-sympower-then-minh}）。
\par
因此在本文体系中必须区分两类“体积”：
\begin{itemize}
  \item \textbf{编译体积 \(\log|N_q|\)}：由对称幂商核/占据数空间决定，属于实现后端的组合体积；
  \item \textbf{可见维度 \(d_q\)}：由 Hankel 秩/最小 realization 决定，才是碰撞序列语义在边界上真正可审计的维数。
\end{itemize}
二者差额是由对称与代数约束诱导的结构冗余量；它不是“优化技巧”，而是可复核的制度性压缩现象（例如由 Hankel null 证书给出）。
\end{remark}

\begin{theorem}[碰撞矩谱的常数证书与 prime-register 可逆化的超线性证书之间的严格分离]\label{thm:pom-macro-moment-vs-micro-prime-register-certificate-separation}
\leanverified{paper\_pom\_macro\_moment\_vs\_micro\_prime\_register\_certificate\_separation}
固定任意整数 \(q\ge 2\)。
设碰撞矩序列 \((S_q(m))_{m\ge 0}\) 在整数意义下满足一个固定阶的最小整系数递推，其首一最小湮灭多项式记为 \(P_q(\lambda)\in\ZZ[\lambda]\)（参见命题 \ref{prop:pom-moment-minrecurrence-hankel} 及与之配套的 Hankel-null 证书口径，定理 \ref{thm:pom-hankel-syndrome-module-kernel-equals-multiples}）。
则以下两类证书资源在尺度 \(m\to\infty\) 下存在严格分离：
\begin{enumerate}
  \item \textbf{矩谱侧（Hankel 综合征/递推证书）}：多项式 \(P_q\)（或其等价的 Hankel 综合征规范证书，例如推论 \ref{cor:pom-hankel-syndrome-hnf-model-free} 的列 Hermite 正规形）为一个与 \(m\) 无关的常数量对象，足以离线核验并生成全部 \(\{S_q(m)\}_{m\ge 0}\)；
  \item \textbf{prime-register 侧（逐步可逆化外置证书）}：对折叠系统 \(\Fold_m\) 的任意确定性可逆外置策略，其 Gödel 整数化证书的位长度满足超线性下界
  \[
  \EE[\log_2 G_{\mathsf S}(\omega)]\ \ge\ c\,m\log m-O(m)
  \qquad(\omega\sim\mathrm{Unif}(\Omega_{m+1}))
  \]
  （其中常数 \(c>0\) 可取为定理 \ref{thm:conclusion-prime-integerization-superlinear-bitlength} 所给值）。
\end{enumerate}
因此，对固定 \(q\) 而言，矩谱侧证书规模可取为 \(O(1)\)（与 \(m\) 无关），而 prime-register 可逆化证书规模必为 \(\Omega(m\log m)\)；二者之比随 \(m\to\infty\) 发散。
\end{theorem}
